% - GitHub API
%   - REST API
%     - Základní vlastnosti: HTTP, JSON, verzování
%     - Autentizace: Personal Access Token (PAT), GitHub App, OAuth
%     - Příklady endpointů relevantních pro monitoring (issues, PRs, commits, labels...)
%   - GraphQL API
%     - Motivace: REST vracelo příliš mnoho/málo dat, GraphQL umožňuje dotaz na přesně požadovaná pole
%     - Základní syntaxe dotazu, ukázka
%     - Výhody pro monitoring velkých repozitářů (méně requestů, přesnější data)
%     - Omezení API (rychlost)
%   - Rate limiting
%     - REST: 5 000 requestů/hod pro autentizované uživatele
%     - GraphQL: bodový systém (cost/complexity)
%     - Hlavičky X-RateLimit-*, strategie pro vyhnutí se limitům (throttling, caching)
%   - Webhooky
%     - Push notifikace při události (PR otevřen/zavřen, push, label přidán...)
%     - Alternativa k polling přístupu, výhody a nevýhody
%     - Typy eventů relevantních pro monitoring

\section{GitHub API}
\label{sec:github_api}
Github má k dispozici dva typy API pro přístup jak k datům tak pro kompletní interakci s platformou github. Github api umožnuje uživatelům automatizovat procesy a integrovat své aplikace s platformou github.
V této sekci jsouposány základní vlastnosti a možnosti obou API, stejně jako omezení, která je třeba mít na paměti při jejich využívání.

\subsection{REST API}
\label{subsec:github_api_rest}
REST API umožnuje přístup k datům a funkcím GitHubu prostřednictvím HTTP požadavků. Každý požadavek na rest api obsajuje http metodu a cestu. Nicméně záleží na konkrétním endpointu jaké parametry musí být specifikovány. například další hlavičky, informaci o autentizaci, query parameters nebo tělo požadavku
Využívá základní HTTP metody (GET, POST, PUT, DELETE) pro interakci s různými zdroji.
\begin{itemize}
  \item GET pro získání dat (např. seznam issues, detail PR)
  \item POST pro vytvoření nových zdrojů (např. vytvoření issue)
  \item PUT pro aktualizaci existujících zdrojů (např. přidání labelu k issue)
  \item DELETE pro odstranění zdrojů (např. smazání komentáře)
  \item PATCH pro aktualizaci properties (např. změna stavu issue)
\end{itemize}
Každý endpoint je definován cestou. Ta může obsahovat jak statické části tak dynamické části. například endpoint pro získání issues od repozitáře je \code{GET /repos/{owner}/{repo}/issues}, kde \code{owner} a \code{repo} jsou dynamické části, které musí být nahrazeny konkrétními hodnotami.

Data jsou věwwtšinou zprostředkována za pomocí JSON formátu což zajištuje dobrou čitelnot jak pro lidi tak i stroje. Ovšem může se použít i jiný typ Media Type, například diff, patch, raw, full a další. Pro použití těchto typů se používá \code{application/vnd.github.PARAM+json} kde \code{PARAM} je nahrazeno konkrétním typem média. Například pro získání dat ve formátu diff se použije \code{application/vnd.github.diff+json}.

github rest api je verzované a využívá kalendářní verzování. tato hlavička není vyžadována ale je silně doporučována. bez specifikace této hlavičky se automaticky dotazuje na verzi specifikovanou v dokumentaci (2022-11-28). Tedy dotazy na kterou verzi api se posílají jsou v hlavičce požadavku ve formátu \code{--header "X-GitHub-Api-Version:2026-03-10"}.tedy budoucí nová verze nezmění způsob fungování starších verzí. github také poskytuje seznam zpětně nekompatibilních změn mezi každou verzí,  návody a dokumentaci pro migrování na novější verze api.

Ikdyž Některé endpointy jsou veřejné i bez autentizace, většina z nich vyžaduje autentizaci pomocí autentizačního tokenu v hlavičce. Existuje více možností jak získat token, například pomocí PAT (personal access token), github app nebo vestavěný \code{GITHUB\_TOKEN} v rámci github actions. Přihlašovánií za pomocí PAT podpruje jemně odstupnovaný přítup oprávnění, který je upřednosňován oproti klasickému kde tak jemně odstupnovaný přístup oprávnění není možný.

\subsection{GraphQL API}
\label{subsec:github_api_graphql}
Github poskytuje i novější GraphQL API, které umožnuje kientům dotazovat se pouze přesně na to co vyžadují. jedná se o jiný způsob přístupu k datům. tento přístup umožnuje vysokou flexibilitu a menší přenosy dat. Rest api vrací předem definované objekty z mnoha endpointů zatímco u graphql je pouze jeden endpoint který umožnuje dotázat se přesně na ta data která jsou potřeba.

Jedny z hlavních problémů rest api byl buď nadměrné stahování dat (over-fetching) a nnebo nedostatečné stahování dat (under-fetching).

\begin{description}
  \item[Over-fetching] nastává když klientská aplikace potřebuje získat například pouze dvě vlastnosti od objektu ale v rest api se vrací celýobjekt se svými mnoha vlastnostmi. To vede k neefektivnímu využidí dat
  \item[under-fetching] nastává když
\end{description}
